\documentclass[aspectratio=169,11pt]{beamer}

\usetheme{Madrid}
\usecolortheme{default}
\usefonttheme{professionalfonts}
\setbeamertemplate{navigation symbols}{}
\setbeamertemplate{footline}[frame number]

\usepackage{amsmath, amssymb, booktabs}

\title{Hiring, Promotion, Pay-Setting, And Firm-Side Gender Gaps}
\subtitle{Gender Study --- Week 4}
\author{Teaching deck draft}
\date{}

\begin{document}

\begin{frame}
  \titlepage
\end{frame}

\begin{frame}{Central question}
\begin{itemize}
    \item How do firms and internal labor markets produce gender differences in pay, promotion, authority, and retention?
    \item Firms are labor-market institutions: they allocate jobs, authority, rents, schedules, metrics, promotion chances, and wage premiums.
    \item Week 4 discipline: separate worker sorting from within-firm behavior and pay-setting rules.
\end{itemize}
\end{frame}

\begin{frame}{Week 3 to Week 4 bridge}
\begin{itemize}
    \item Week 3: education, aspirations, skills, and early occupations shape who reaches which firms and career tracks.
    \item Week 4: once workers arrive, firms screen, supervise, evaluate, pay, promote, and retain.
    \item The same initial sorting can lead to different careers under different managers, tasks, pay rules, and promotion systems.
\end{itemize}
\end{frame}

\begin{frame}{Three firm-side channels}
\small
\begin{tabular}{p{0.25\linewidth} p{0.55\linewidth}}
\toprule
Channel & Core object \\
\midrule
Sorting into firms and jobs & exposure to high-premium firms, ladders, teams, schedules, conditions \\
Within-firm personnel processes & hiring, referrals, supervision, evaluation, task allocation, promotion, networks \\
Pay-setting and bargaining & wage premiums, salary requests, negotiation, transparency, pay bands, compliance \\
\bottomrule
\end{tabular}
\end{frame}

\begin{frame}{Conceptual map}
\[
\begin{aligned}
G_Y
&=
\Delta^{sort}_{Y}
+ \Delta^{personnel}_{Y}
+ \Delta^{pay}_{Y}
+ \Delta^{retention}_{Y}
+ \Delta^{interaction}_{Y}
\end{aligned}
\]
\begin{itemize}
    \item \(Y\): wages, earnings, authority, promotion, hours, job quality, or retention.
    \item Sorting is location in firms and jobs; personnel is treatment after entry.
    \item Pay-setting and retention can reinforce both sorting and promotion.
\end{itemize}
\end{frame}

\begin{frame}{Sorting is not firm behavior}
\begin{itemize}
    \item Workers may enter different employers, establishments, teams, schedules, occupations, and ladders.
    \item High-premium firms can differ in rents, promotion chances, flexibility, safety, and authority.
    \item A between-firm gap does not by itself prove within-firm differential treatment.
    \item Matched employer-employee data help separate firm exposure from within-firm pay.
\end{itemize}
\end{frame}

\begin{frame}{Hiring and screening}
\[
H_{ij}=\mathbf{1}\{S_{ij}(X_i,A_i,G_i,R_i)\ge c_j\}
\]
\begin{itemize}
    \item Hiring designs change what the screen observes or compare otherwise similar applications.
    \item Blind-screen and audit designs are clean for entry-stage advancement.
    \item The observed margin is callback, advancement, interview, hire, or starting job quality.
    \item Do not turn a hiring estimate into a promotion or retention estimate.
\end{itemize}
\end{frame}

\begin{frame}{Blind auditions as an entry design}
\begin{itemize}
    \item The classic object is a procedural change in screening.
    \item Identifying variation: blind versus nonblind audition conditions.
    \item Observed margin: advancement and hiring at entry.
    \item The economic lesson: information design can change who passes the screen.
\end{itemize}
\end{frame}

\begin{frame}{Promotion and authority pipeline}
\[
P_{i,t+1}=f(P_{it},Q_{it},M_{it},N_{it},T_{it},E_{it})
\]
\begin{itemize}
    \item \(Q_{it}\): measured performance.
    \item \(M_{it}\), \(N_{it}\): manager access, mentoring, sponsorship, network capital.
    \item \(T_{it}\): tasks that create visibility or promotable output.
    \item \(E_{it}\): evaluation and potential ratings.
\end{itemize}
\end{frame}

\begin{frame}{Potential ratings and leadership ladders}
\begin{itemize}
    \item Firms promote on expected future leadership, not only current output.
    \item Subjective potential can encode information, bias, risk assessment, or retention strategy.
    \item Personnel records matter because they can observe performance, ratings, levels, and promotions together.
    \item The margin is authority allocation, not just wage growth.
\end{itemize}
\end{frame}

\begin{frame}{Networks, mentoring, and old-boy channels}
\begin{itemize}
    \item Internal careers often depend on social access to managers and sponsors.
    \item Informal interaction can affect information, visibility, advocacy, and nominations.
    \item Strong designs use manager rotation, reassignment, overlap, or shift timing.
    \item Survey-linked administrative data can observe hidden channels and personnel outcomes.
\end{itemize}
\end{frame}

\begin{frame}{Task allocation and low-promotability work}
\begin{itemize}
    \item Firms allocate tasks as well as wages.
    \item Some tasks build revenue, visibility, skill, client access, or authority.
    \item Low-promotability tasks are necessary but weakly rewarded.
    \item If women receive or accept more of these tasks, promotion gaps can grow even when formal pay rules look neutral.
\end{itemize}
\end{frame}

\begin{frame}{Firm-specific pay premiums}
\[
\begin{aligned}
w_{ijt}
&=\alpha_i+\psi_j+X_{it}'\beta+\theta_j Female_i+\varepsilon_{ijt}
\end{aligned}
\]
\begin{itemize}
    \item \(\psi_j\): firm wage premium.
    \item Sorting channel: who works at high-\(\psi_j\) firms?
    \item Within-firm channel: does the firm share rents differently?
    \item Matched employer-employee decompositions are strong for this separation.
\end{itemize}
\end{frame}

\begin{frame}{Bargaining and pay-setting}
\begin{itemize}
    \item Firms choose whether wages come from posted ranges, salary histories, requests, negotiation, pay bands, or manager discretion.
    \item Salary requests and outside offers can become pay-setting inputs.
    \item The same request gap can matter more in firms with more discretionary wage-setting.
    \item Ask which rule transformed information into pay.
\end{itemize}
\end{frame}

\begin{frame}{Transparency and equal-pay frontier}
\begin{itemize}
    \item Pay transparency changes information available to workers, managers, and regulators.
    \item Equal-pay rules target comparable jobs and internal pay consistency.
    \item Firm responses can include wage compression, raises, job reclassification, changed posting behavior, or changed sorting.
    \item Policy evidence must name the treated unit, comparison group, and margin moved.
\end{itemize}
\end{frame}

\begin{frame}{Retention and workplace conditions}
\begin{itemize}
    \item Exit is a firm-side margin when workplace conditions affect who stays.
    \item Harassment, weak protection, unpredictable schedules, exclusion, and poor supervision can change the cost of remaining.
    \item Administrative quits can hide job-quality losses and selection out of high-wage paths.
    \item Retention connects firm-side inequality to safety and mobility later in the course.
\end{itemize}
\end{frame}

\begin{frame}{Designs and margins}
\scriptsize
\setlength{\tabcolsep}{3pt}
\begin{tabular}{p{0.23\linewidth} p{0.25\linewidth} p{0.27\linewidth}}
\toprule
Design & Identifying variation & Observed margin \\
\midrule
Audit / blind screen & identity signal or procedure & callback, advancement, hire \\
Personnel records & ratings, levels, performance & promotion, authority, pay growth \\
Manager assignment & rotation, overlap, reassignment & mentoring, visibility, promotion \\
Matched employer-employee & mobility and firm premiums & sorting, rent sharing, pay gap \\
Firm policy reform & treated vs untreated exposure & transparency, compression, retention \\
Survey-linked admin & measured hidden channels & networks, climate, mentoring, exit \\
\bottomrule
\end{tabular}
\end{frame}

\begin{frame}{Week 4 lab}
\begin{itemize}
    \item Reproduce: synthetic manager-access and promotion exercise inspired by Cullen and Perez-Truglia.
    \item Diagnose: sort each output into sorting, personnel, pay-setting, or retention.
    \item Transfer: synthetic blind-screening exercise inspired by Goldin and Rouse.
    \item Optional frontier prompt: pay transparency or subjective potential ratings.
\end{itemize}
\end{frame}

\begin{frame}{Bridge to Weeks 5 and 6}
\begin{itemize}
    \item Week 5: norms and bargaining operate through self-promotion, social access, negotiation, authority, and perceived commitment.
    \item Week 6: policies move margins such as pay rules, disclosure, compliance, leave, quotas, protection, and retention.
    \item The question becomes: which firm margin moved, for whom, and at what cost?
\end{itemize}
\end{frame}

\begin{frame}{Takeaways}
\begin{itemize}
    \item Firms are active allocators of pay, authority, tasks, schedules, and career opportunities.
    \item Separate sorting into firms from within-firm treatment and firm pay-setting.
    \item Hiring, promotion, task allocation, pay, transparency, and retention require different evidence.
    \item Strong empirical work names the identifying variation and the margin observed.
\end{itemize}
\end{frame}

\end{document}
